% =====================================================================
\section{Findings and Results}
% =====================================================================

\subsection{Overall Migration Results}
Over twelve months, three developers used the system to migrate 39
distinct identifiers, producing 595 submitted code changes and
93{,}574 individual edits in total. Table~\ref{tab:aggregate}
summarizes the headline numbers reported in the paper.

\begin{table}[t]
\caption{Aggregate statistics reported for the case study (reproduced from Table~1 of the original paper).}
\label{tab:aggregate}
\begin{tabular}{@{}lr@{}}
\toprule
Metric & Value \\
\midrule
Identifiers migrated & 39 \\
Developers involved & 3 \\
Total code changes submitted & 595 \\
\quad LLM-Only & 214 (35.97\%) \\
\quad LLM-then-Human & 229 (38.48\%) \\
\quad Human-Only & 152 (25.55\%) \\
Distinct reviewers & 306 \\
Distinct teams involved & 149 \\
Total edits (all IDs) & 93{,}574 \\
\quad Attributed to the LLM (LLM$\Delta$) & 64{,}996 (69.46\%) \\
\quad Attributed to humans (Human$\Delta$) & 28{,}578 (30.54\%) \\
\bottomrule
\end{tabular}
\end{table}

In other words, 74.45\% of all submitted code changes contained at
least some LLM-authored content (LLM-Only plus LLM-then-Human), and
the LLM was responsible for roughly two-thirds of all edited
characters across the whole study. The size of individual migrations
varied enormously. Some identifiers required edits spread across
hundreds of files in multiple languages, while others were confined
to a single file.

\subsection{Results Across Individual Migrations}
Table~\ref{tab:perid} reproduces the paper's full per-identifier
breakdown, sorted by descending total edit volume (Total$\Delta$).
The two largest migrations, ID-30 and ID-32, together account for a
disproportionate share of all edits in the study, roughly half of the
total across all 39 identifiers combined, while the smallest, such as
ID-10, ID-24, and ID-39, touched only a handful of files. Larger
identifiers also tended to span more programming languages at once:
ID-32 alone touched five different languages, versus a single
language for most of the smaller identifiers.


\begin{table}[t]
\caption{Per-identifier statistics, sorted by descending Total$\Delta$ (reproduced from Table~2 of the original paper).}
\label{tab:perid}
\footnotesize
\begin{tabular}{@{}lrrrr@{}}
\toprule
ID & Total$\Delta$ & \#Changes & \#Files & \#Langs \\
\midrule
ID-30 & 23{,}800 & 39 & 418 & 2 \\
ID-32 & 22{,}822 & 185 & 392 & 5 \\
ID-20 & 16{,}406 & 28 & 175 & 2 \\
ID-28 & 7{,}059 & 29 & 89 & 1 \\
ID-03 & 4{,}594 & 21 & 350 & 4 \\
ID-37 & 2{,}774 & 25 & 58 & 1 \\
ID-06 & 2{,}458 & 23 & 41 & 2 \\
ID-35 & 2{,}039 & 14 & 38 & 2 \\
ID-19 & 1{,}937 & 17 & 21 & 1 \\
ID-18 & 1{,}889 & 20 & 33 & 1 \\
ID-29 & 1{,}825 & 6 & 38 & 1 \\
ID-04 & 1{,}279 & 10 & 27 & 1 \\
ID-31 & 1{,}223 & 15 & 31 & 1 \\
ID-26 & 1{,}151 & 4 & 11 & 1 \\
ID-34 & 1{,}088 & 4 & 16 & 2 \\
ID-09 & 1{,}038 & 6 & 16 & 1 \\
ID-38 & 844 & 8 & 17 & 2 \\
ID-14 & 802 & 7 & 12 & 2 \\
ID-01 & 798 & 16 & 103 & 4 \\
ID-07 & 673 & 7 & 31 & 1 \\
ID-05 & 584 & 5 & 9 & 1 \\
ID-08 & 578 & 6 & 13 & 2 \\
ID-22 & 571 & 2 & 19 & 1 \\
ID-17 & 538 & 4 & 6 & 1 \\
ID-36 & 386 & 5 & 8 & 2 \\
ID-12 & 374 & 5 & 40 & 4 \\
ID-02 & 367 & 10 & 14 & 1 \\
ID-23 & 326 & 10 & 15 & 1 \\
ID-33 & 299 & 7 & 7 & 2 \\
ID-13 & 270 & 8 & 14 & 2 \\
ID-27 & 221 & 5 & 6 & 1 \\
ID-15 & 221 & 11 & 12 & 2 \\
ID-16 & 156 & 9 & 13 & 2 \\
ID-21 & 146 & 3 & 10 & 1 \\
ID-25 & 128 & 7 & 11 & 1 \\
ID-10 & 98 & 1 & 4 & 1 \\
ID-11 & 95 & 7 & 7 & 2 \\
ID-24 & 26 & 2 & 2 & 1 \\
ID-39 & 23 & 1 & 1 & 1 \\
\bottomrule
\end{tabular}
\end{table}


The LLM$\Delta$ share per identifier ranged from close to 100\% for
some identifiers (e.g., ID-10 and ID-39) down to well under half for
others (e.g., ID-15 at 25.34\%). The paper attributes this spread
mainly to two factors discussed further below: how far along
developers were in ramping up on the tool when a given identifier was
tackled, and how much of that identifier's code lived in a language,
such as Dart, that the underlying model handled less well.

\subsection{Developer Experience}
In post-study interviews, the three developers consistently praised
three aspects of the system. First, \emph{end-to-end automation} of
both change discovery and generation gave them a ``continuously
renewed'' migration backlog rather than a one-time, overwhelming task
list. They estimated a roughly 50\% reduction in total time spent
compared with earlier, fully manual migrations, an estimate the
authors note is directionally consistent with the 74.45\% and 69.46\%
automation figures, even though no precise wall-clock timing data was
collected. Second, \emph{automated validation} meant developers no
longer had to manually kick off and babysit build and test runs
before deciding whether a change was ready for review. Third, the
\emph{flexibility of prompting an LLM in natural language} meant the
same simple instruction worked across several languages and coding
idioms, removing the burden of maintaining bespoke transformation
rules for each one.

\subsection{Limitations of the Automated Workflow}
The system was not without limitations, several of which show up
directly in the statistics. \emph{Non-determinism}: because repeated
LLM calls on identical input can produce different diffs, developers
had mixed reactions, sometimes pleased that a retry eventually
succeeded, sometimes frustrated by not knowing how many retries would
be needed before falling back to manual edits. \emph{Context-window
limits}: very large files sometimes could not be passed to the model
in their entirety, causing automatic edits to fail outright.
\emph{Hallucination}: the model sometimes only reformatted a file,
only added a comment next to an already-correct value, or made
unrelated edits, none of which counted as a genuine fix. Hallucinated outputs
 included formatting-only changes, unnecessary comments, and unrelated edits. 
 However, the Human-Only and Human$\Delta$ proportions should not be interpreted as 
 hallucination rates, since manual work also resulted from context 
 limitations, uneven language support, validation failures, and other 
 workflow constraints. \emph{Uneven language support}: the model
handled Java, C++, and Python well but had comparatively weak Dart
support, so identifiers with heavy Dart usage ended up with lower
LLM$\Delta$ shares. \emph{Ramp-up effects}: early in the study,
developers were still learning the tool and tackled small, localized
identifiers largely by hand; as they gained experience they moved on
to larger, more automation-heavy migrations, which the authors say
explains much of the variance seen across individual identifiers.

\subsection{Pre-Existing Failures and Production Rollouts}
Two further sources of friction were largely outside the automated
pipeline's control. First, every code change at Google is expected to
pass the relevant regression test suite before submission; if that
suite already had unrelated, pre-existing failures, the migration's
own validation step would also report a failure even when the LLM's
edit was correct, and developers had to separately track down the
code owners to fix the unrelated failure first. Second, some tests
store their expected output as a serialized ``golden'' file,
regenerated by a specialized tool whenever the code under test
changes. Because the automated pipeline had no way to invoke these
golden-regeneration tools, any test relying on stored goldens would
fail the pipeline's test-validation step regardless of whether the
underlying code change was correct, again pushing that file into the
manual, Human-Only bucket.

Most of the 39 identifiers were small and localized enough that, once
submitted, no further coordination was needed. The two largest and
most widely used identifiers, ID-30 and ID-32, were the exception:
because their blast radius potentially touched many downstream teams,
the resulting changes were rolled out to production gradually, with
affected teams watching for regressions, rather than shipped all at
once. The automated pipeline offered no direct support for this
staged-rollout process; it remained a manual, judgment-driven step
carried out by the developers and code owners together.